\chapter{Random Variables and Distributions}

A messaging service receives an average of three requests per second at peak hours.
That number is printed on the dashboard, and it does not answer the pressing question: how
often does a single second hold eight requests or more, so that the queue overflows? Two systems can have exactly the same average and behave
very differently: one almost never overflows, the other several times a day.
The average alone is not enough, and not because the counting is insufficiently careful.

Chapter 17 gave us a language for events: the probability of one outcome, the probability of a union,
conditional probability. That language stops working as soon as what we measure is a number.
``The number of requests in one second'' is not one event but a family of
infinitely many events, one for each possible value. For
counts, the probability of each value can still be added up one by one. For a quantity that
is continuous, that way collapses: a waiting time can be $0.317$ seconds or
$0.3170001$ seconds, and there is no way to state ``the probability of exactly $0.317$ seconds''
without assigning weight zero to every value.

So two things are needed: a way to move probability onto the number line, and a way
to replace summation by integration when the values are continuous.
Once both exist, the average, the spread, and the shape of the tail can be computed.

\begin{untukanak}
Imagine a bakery. Every hour, the seller writes down how many loaves were
sold: sometimes 3, sometimes 11, sometimes none at all. Those records are
whole numbers, and cannot take the value one half. If we want to know the probability
of selling exactly 7 loaves, we can compute it: count how many hours sold 7 loaves,
divide by the number of hours recorded. All the probabilities can be added one by one,
because the options are few and can be listed.

Now weigh one loaf on a digital scale. The result is not a whole number:
250.3 grams, or 249.8 grams. However many decimal places we take, there is always
another possibility closer still. Hence the question ``what is the probability that the weight is
exactly 250.0000 grams'' is useless, whereas ``what is the probability that the weight lies between 249
and 251 grams'' is very useful. Probability here belongs not to a single point but to
an interval.

In this chapter we will learn why those two ways differ, and why both can still
give a sensible average.
\end{untukanak}

\section{From Outcomes to Numbers}
\tingkat{basic}

A measuring instrument does not record events; it records numbers. The bridge is a function
that takes each outcome of the experiment and assigns one real number. The name
``random variable'' is misleading: it is not a variable, and it is not random. The randomness
comes from the probability on the outcome space, whereas the function is ordinary.

\begin{definisi}[Random variable]\index{random variable}
Let $\Omega$ be the outcome space and $\mathcal{F}$ a collection of events. The mapping
$X:\Omega\to\mathbb{R}$ is called a \emph{random variable} if for every $x\in\mathbb{R}$
the set
\[
\{X\le x\}=\{\omega\in\Omega : X(\omega)\le x\}
\]
is an event, that is, a member of $\mathcal{F}$.
\end{definisi}

That last condition guarantees that $P(X\le x)$ has a meaning. For almost every function that
appears in practice it holds automatically, and it fails only for set indicators that
are not measurable.

As an example, toss three fair coins and let $X$ be the number of heads that appear.
Its values are $0,1,2,3$ with probabilities $1/8,3/8,3/8,1/8$. Note that $X$ maps
eight outcomes onto four numbers, so several outcomes share the same value. That is
why the probability of a value is the probability of a union of events.

We need one tool that answers ``how far have the values moved already'' for all
thresholds at once, because the threshold is chosen by the reader, not by the author of the book.

\begin{definisi}[Cumulative distribution function]\index{cumulative distribution function}
The cumulative distribution function of a random variable $X$ is
$F_X(x)=P(X\le x)$ for $x\in\mathbb{R}$.
\end{definisi}

\begin{proposisi}[Properties of the distribution function]
For every random variable $X$ the following hold:
\begin{enumerate}
  \item $F_X$ is nondecreasing: if $a\le b$ then $F_X(a)\le F_X(b)$;
  \item $\lim_{x\to-\infty}F_X(x)=0$ and $\lim_{x\to+\infty}F_X(x)=1$;
  \item $F_X$ is right-continuous, that is, $F_X(x)=F_X(x^+)$ for every $x$;
  \item $P(a<X\le b)=F_X(b)-F_X(a)$ for every $a<b$.
\end{enumerate}
\end{proposisi}

\begin{proof}
For (1), the event $\{X\le a\}$ is contained in $\{X\le b\}$ when $a\le b$, and
probability is monotone with respect to inclusion.

For (2), let $A_n=\{X\le -n\}$. These events are decreasing, and
$\bigcap_n A_n=\varnothing$ because every outcome yields a finite value. Continuity of
probability from above, which follows from countable additivity, gives $F_X(-n)\to0$. The
opposite direction uses $B_n=\{X\le n\}$, which increase to $\Omega$, so that
$F_X(n)\to1$.

For (3), take a decreasing sequence $y_n\downarrow x$ and $C_n=\{X\le y_n\}$. Then
$\bigcap_n C_n=\{X\le x\}$, hence $F_X(y_n)\to F_X(x)$ by continuity from above.
Since the sequence is arbitrary, the right-hand limit equals the value at $x$.

For (4), the event $\{X\le b\}$ splits into $\{X\le a\}$ and $\{a<X\le b\}$,
so that $F_X(b)=F_X(a)+P(a<X\le b)$.
\end{proof}

For two tosses of a fair coin with $X$ the number of heads,
\[
F_X(x)=\begin{cases}
0, & x<0,\\
1/4, & 0\le x<1,\\
3/4, & 1\le x<2,\\
1, & x\ge 2.
\end{cases}
\]
This function is defined on the whole number line, even though $X$ takes only the values $0,1,2$.
That is why $F_X(1.4)=3/4$: the question is not ``is $X$ equal to $1.4$'',
but ``is $X$ at most $1.4$''. Every jump of the graph is the probability
mass at a single value. Figure~\ref{fig:tangga-cdf} draws that staircase: three horizontal
steps, three vertical jumps, and the height of each jump equals the probability
mass at the point of the jump.

\begin{figure}[htbp]
\centering
\begin{tikzpicture}[font=\footnotesize,>=Stealth,x=1.75cm,y=4.2cm]
  \foreach \y/\lab in {0.25/{1/4},0.75/{3/4},1/{1}}{
    \draw[layergrey!35,dashed] (-1.35,\y) -- (3.55,\y);
    \node[left,font=\scriptsize] at (0,\y) {$\lab$};
  }
  \draw[->,layergrey] (-1.35,0) -- (3.6,0) node[below] {$x$};
  \draw[->,layergrey] (0,-0.1) -- (0,1.18) node[left] {$F_X(x)$};
  \draw[very thick,layerindigo] (-1.25,0) -- (0,0);
  \draw[very thick,layerindigo] (0,0.25) -- (1,0.25);
  \draw[very thick,layerindigo] (1,0.75) -- (2,0.75);
  \draw[very thick,layerindigo] (2,1) -- (3.5,1);
  \draw[very thick,layerplum] (0,0) -- (0,0.25);
  \draw[very thick,layerplum] (1,0.25) -- (1,0.75);
  \draw[very thick,layerplum] (2,0.75) -- (2,1);
  \fill[layerindigo] (0,0.25) circle (1.8pt);
  \fill[layerindigo] (1,0.75) circle (1.8pt);
  \fill[layerindigo] (2,1) circle (1.8pt);
  \draw[layerindigo,fill=white,thick] (1,0.25) circle (1.8pt);
  \draw[layerindigo,fill=white,thick] (2,0.75) circle (1.8pt);
  \node[below,font=\scriptsize] at (0,-0.02) {$0$};
  \node[below,font=\scriptsize] at (1,-0.02) {$1$};
  \node[below,font=\scriptsize] at (2,-0.02) {$2$};
  \node[right,font=\scriptsize,text=layerplum] at (1.06,0.5) {jump $1/2$};
  \node[right,font=\scriptsize,text=layerplum] at (0.06,0.12) {jump $1/4$};
\end{tikzpicture}
\caption{The cumulative distribution function of the number of heads in two tosses of a coin. The graph is
a flat staircase between the possible values, then jumps exactly at those values. The jumps of
$1/4$ at $x=0$, $1/2$ at $x=1$, and $1/4$ at $x=2$ are the probability masses $p_X(0)$,
$p_X(1)$, and $p_X(2)$.
The filled dot at the left end of every step marks that $F_X$ is right-continuous, so that
$F_X(1)=3/4$ and not $1/4$.}
\label{fig:tangga-cdf}
\end{figure}

\kamus{random variable}{variabel acak}{a function that attaches one number to each
outcome of the experiment, like a clerk who records one number for every event}

\kamus{cumulative distribution function}{fungsi distribusi kumulatif}{a rising
staircase: its height at the point $x$ states how much probability has already
accumulated up to that threshold}

\section{Mass Functions and Density Functions}
\tingkat{basic}

There are two ways of weighting the values of a random variable. For a variable whose values can
be listed, the weight attaches to points; for one whose values fill an interval, the weight attaches
to length. The most common mistake in this topic is to use the first way to
read the outcome of the second.

\begin{definisi}[Discrete variables and probability mass function]\index{probability mass function}
$X$ is called \emph{discrete} if its set of values is countable, say
$\{x_1,x_2,\dots\}$. The \emph{probability mass function} of $X$ is $p_X(x)=P(X=x)$.
This function satisfies $p_X(x)\ge0$, $p_X(x)\le1$, and $\sum_i p_X(x_i)=1$, as well as
$F_X(x)=\sum_{x_i\le x}p_X(x_i)$.
\end{definisi}

For $X$ binomial with $n=2$ and $p=1/2$, the probability mass is $1/4,1/2,1/4$ at
$x=0,1,2$ and zero elsewhere. Probability mass is never greater than one,
because it is a probability after all.

\begin{definisi}[Continuous variables and probability density function]\index{density function}
$X$ is called \emph{continuous} if there exists a function $f_X\ge0$ such that
\[
P(a\le X\le b)=\int_a^b f_X(x)\,dx
\]
for every $a\le b$, with $\int_{-\infty}^{\infty}f_X(x)\,dx=1$. The function $f_X$
is called the \emph{probability density function}.
\end{definisi}

Density is not probability, and the difference is not a matter of terminology. Take $X$ with
density $f_X(x)=5$ on $[0,1/5]$ and zero outside it. Then $f_X(0.1)=5$, a number
that is impossible for a probability, whereas
$P(X\le0.1)=\int_0^{0.1}5\,dx=0.5$. Density may also be unbounded:
$f_X(x)=\tfrac23x^{-1/3}$ on $(0,1)$ integrates to one even though it blows up near zero.
The units differ: probability per unit of length, just as population density is people
per square kilometer. Figure~\ref{fig:batang-dan-kurva} puts
the two ways side by side: the bars on the left are added by height, the triangle on the
right by area, and both total one.

\begin{figure}[htbp]
\centering
\begin{tikzpicture}[font=\footnotesize,>=Stealth]
  \begin{scope}[x=1.05cm,y=5.4cm]
    \draw[->,layergrey] (0,0) -- (5.4,0) node[right] {$x$};
    \draw[->,layergrey] (0,0) -- (0,0.62) node[left] {$p_X(x)$};
    \foreach \x/\p in {1/0.1,2/0.2,3/0.3,4/0.4}{
      \fill[layerorange!40] (\x-0.3,0) rectangle (\x+0.3,\p);
      \draw[layerbrown,thick] (\x-0.3,0) rectangle (\x+0.3,\p);
      \node[below,font=\scriptsize] at (\x,0) {$\x$};
    }
    \node[font=\scriptsize,text=layerbrown,align=left,anchor=west] at (0.45,0.44)
      {four bars,\\ total height $=1$};
  \end{scope}
  \begin{scope}[xshift=7.7cm,x=1.05cm,y=5.4cm]
    \draw[->,layergrey] (0,0) -- (5.4,0) node[right] {$x$};
    \draw[->,layergrey] (0,0) -- (0,0.62) node[left] {$f_X(x)$};
    \fill[layerteal!25] (0.4,0) -- (2.4,0.5) -- (4.4,0) -- cycle;
    \draw[layerteal,thick] (0.4,0) -- (2.4,0.5) -- (4.4,0) -- cycle;
    \draw[layergrey,dashed] (2.4,0) -- (2.4,0.5);
    \node[above,font=\scriptsize] at (2.4,0.5) {$0.5$};
    \node[below,font=\scriptsize] at (0.4,0) {$0$};
    \node[below,font=\scriptsize] at (4.4,0) {$4$};
    \node[font=\scriptsize,text=layerteal,align=center] at (2.4,0.14)
      {area $=\tfrac12\cdot4\cdot0.5=1$};
  \end{scope}
\end{tikzpicture}
\caption{Two ways of weighting that both total one. Left: four bars
of height $0.1$, $0.2$, $0.3$, $0.4$ at the values $1,2,3,4$, added by
height. Right: a triangle-shaped density with peak $0.5$ on $[0,4]$,
added by area, and $\tfrac12\cdot4\cdot0.5=1$. The height $0.5$ in the right-hand panel
is not a probability of anything; what has probability is the area of the slice beneath it.}
\label{fig:batang-dan-kurva}
\end{figure}

\begin{bacakeras}
Read aloud in pairs, one person taking the Height voice and one the Weight voice. Start
slowly, then speed up.

Height: I am the height of the curve. I tower at one point, and everyone points at me.
Weight: I am the weight under the curve. I pile up along the interval, and everyone weighs me.
Height: My number may be five. May be ten. May be unbounded.
Weight: My number never exceeds one, because I am the whole of probability.
Height: Without me, you do not know where the most frequent value lies.
Weight: Without me, you do not know whether that peak contains anything.
Together: Height points to the place. Weight weighs the content.

Repeat together three times, more slowly each time: probability is weight, not height.
\end{bacakeras}

\begin{teorema}\index{density}
If $X$ is continuous with density $f_X$, then $F_X(x)=\int_{-\infty}^{x}f_X(t)\,dt$,
and $f_X(x)=F_X'(x)$ at every point $x$ where $F_X$ is differentiable.
\end{teorema}

\begin{proof}
The first equation is the definition of density with $a=-\infty$ and $b=x$. The second
equation is the first part of the fundamental theorem of calculus: the derivative of the accumulation
function equals $f_X$ at every point where $f_X$ is continuous.
\end{proof}

The question that most often causes confusion: if $X$ is continuous then
$P(X=x)=0$ for every $x$, because $\int_x^x f_X=0$. Does that mean that the event $X=x$
is impossible? No. The set $\{\omega:X(\omega)=x\}$ still exists and can still contain
outcomes; what is zero is its probability, not its existence.

The difficulty appears when we try to add up all the points. The interval $[0,1]$ is an uncountable
union of points of probability zero, yet $P(0\le X\le1)=1$. There is no
contradiction: the third axiom of Chapter 17 guarantees additivity only for
\emph{countable} unions, and the integral measures length, it does not count points.

There is a practical layer. NumPy's built-in generator forms values
as multiples of $2^{-53}$, so there are only $2^{53}$ possible outputs on $[0,1)$, and
the probability that one particular value appears in $10^7$ draws is about
$10^7/2^{53}=1.1\times10^{-9}$. On a computer ``probability zero'' is really
$2^{-53}\approx1.1\times10^{-16}$; the zero we write is the idealization.

\begin{contoh}
Let $X$ be uniform on $(0,1)$, with $f_X(x)=1$ for $0<x<1$. Then
\[
P(0.49<X<0.51)=\int_{0.49}^{0.51}1\,dx=0.02,
\qquad P(X=0.5)=0 .
\]
The density around $0.5$ is as large as the density around $0.01$; what distinguishes them
is only the length of the interval. That is why the right question for a continuous variable
is always an interval.
\end{contoh}

\begin{tangan}
Materials: a tape measure, scissors, card stock, and ten blank cards.

Measure the heights of ten friends to the nearest centimeter, then write one height on each
card. Make five boxes on the floor from strips of card stock five centimeters wide:
$150$--$154$, $155$--$159$, up to $170$--$174$. Put each card in the box that
matches, stacking them upward. What we must see: those stacks have a shape, and the number of
cards in all the boxes is still ten.

Cut the card stock again into boxes ten centimeters wide, arrange the cards again,
and see what changes. The stacks become almost half as high, but not a single card
is lost. The width of the box changes the height of the stack; the number of cards does not
change. That is the difference between height and weight, and that is why a density must be multiplied
by the width of the interval before it may be called a probability.
\end{tangan}

\begin{untukanak}
Imagine a football field full of people. We want to know where the people are
most numerous. We cannot answer that by asking ``how many people are exactly at
this point''. The answer is always zero, because a point has no area. What can
be answered is ``how many people are there per square meter''. That number is called density,
and it differs from place to place: dense at the entrance, sparse in the corner of the field.

High density does not mean that there are many people at one point there. High
density means that around that place, within a small region, people are packed together. If we
multiply the density by the area of the region, only then does the number of people come out. The same
holds for probability: density times the length of the interval gives probability, and density
itself is not probability.

This is why two curves of very different shape can give the same total,
namely one. What matters is not the height of the curve at one point, but the area beneath it.
In this chapter we will learn why a question about points must be replaced by a question
about intervals.
\end{untukanak}

\section{Expectation and Variance}
\tingkat{intermediate}

If the whole distribution has to be compressed into a single number, the number chosen
is its balance point: the location of the center of mass when probability is treated as weight.

\begin{definisi}[Expectation]\index{expectation}
The expectation of a random variable $X$ is
\[
\mathbb{E}[X]=\sum_i x_i\,p_X(x_i)
\qquad\text{or}\qquad
\mathbb{E}[X]=\int_{-\infty}^{\infty}x\,f_X(x)\,dx,
\]
depending on whether $X$ is discrete or continuous. The expectation exists only if
$\mathbb{E}[|X|]<\infty$; otherwise the expectation of $X$ is said not to exist.
\end{definisi}

That absolute-value condition is not excessive caution. For the Cauchy distribution with
density $f(x)=1/(\pi(1+x^2))$, $\int|x|f(x)dx$ diverges, so that $\mathbb{E}[X]$
does not exist even though the distribution is symmetric about zero; the sample mean does not
concentrate anywhere, however much data we add. Chapter 19 uses this example
as a boundary of the law of large numbers.

As short computations, for $X$ Bernoulli with $P(X=1)=p$,
$\mathbb{E}[X]=0\cdot(1-p)+1\cdot p=p$. For $X$ uniform on $(0,1)$,
$\mathbb{E}[X]=\int_0^1x\,dx=1/2$.

What is often needed is not the expectation of $X$ but the expectation of a function of $X$, for example
$\mathbb{E}[X^2]$. The direct route is to find the density of $X^2$ first, then
integrate. There is a much shorter route.

\begin{teorema}[Expectation of a function of a random variable]
For a function $r$ and a random variable $X$,
\[
\mathbb{E}[r(X)]=\sum_i r(x_i)\,p_X(x_i)
\qquad\text{or}\qquad
\mathbb{E}[r(X)]=\int_{-\infty}^{\infty}r(x)\,f_X(x)\,dx .
\]
\end{teorema}

\begin{proof}
For the discrete case, group the values according to $y=r(x)$:
\[
\sum_i r(x_i)p_X(x_i)=\sum_y y\sum_{i:\,r(x_i)=y}p_X(x_i)=\sum_y y\,P(r(X)=y),
\]
and the last expression is the definition of $\mathbb{E}[r(X)]$ because $r(X)$ is discrete with
probability mass $P(r(X)=y)$. For the continuous case, the proof is the change-of-variables
formula for integrals in Chapter 15; the same grouping applies with integrals
replacing sums.
\end{proof}

\begin{proposisi}[Linearity of expectation]
For random variables $X,Y$ and constants $a,b$,
$\mathbb{E}[aX+bY]=a\,\mathbb{E}[X]+b\,\mathbb{E}[Y]$.
\end{proposisi}

\begin{proof}
For the discrete case with joint mass $p(x,y)$,
\[
\mathbb{E}[aX+bY]=\sum_{x,y}(ax+by)p(x,y)
=a\sum_x x\sum_y p(x,y)+b\sum_y y\sum_x p(x,y)=a\,\mathbb{E}[X]+b\,\mathbb{E}[Y].
\]
The continuous case is identical with integrals replacing sums.
\end{proof}

What is not used here is independence: linearity holds for any variables,
including strongly dependent ones. Hence it is safe for sums of indicators, but it may not
be used for products: $\mathbb{E}[XY]=\mathbb{E}[X]\mathbb{E}[Y]$ requires
independence.

Expectation gives the center, but it does not give the width; the most widely used measure of width
is the average squared distance to the center.

\begin{definisi}[Variance]\index{variance}
The variance of $X$ with $\mu=\mathbb{E}[X]$ is
$\operatorname{Var}(X)=\mathbb{E}\bigl[(X-\mu)^2\bigr]$, and its standard deviation is
$\sigma_X=\sqrt{\operatorname{Var}(X)}$.
\end{definisi}

\begin{proposisi}
If the variance exists, then $\operatorname{Var}(X)=\mathbb{E}[X^2]-\mu^2$ and
$\operatorname{Var}(aX+b)=a^2\operatorname{Var}(X)$.
\end{proposisi}

\begin{proof}
Expand the square and use linearity:
\[
\mathbb{E}[(X-\mu)^2]=\mathbb{E}[X^2]-2\mu\,\mathbb{E}[X]+\mu^2
=\mathbb{E}[X^2]-\mu^2 .
\]
For the second property, $\mathbb{E}[aX+b]=a\mu+b$, so that
$\operatorname{Var}(aX+b)=\mathbb{E}[(aX+b-a\mu-b)^2]=a^2\mathbb{E}[(X-\mu)^2]$.
\end{proof}

For $X$ Bernoulli, $\mathbb{E}[X^2]=p$ because $X^2=X$ for the values $0$ and $1$,
so that $\operatorname{Var}(X)=p-p^2=p(1-p)$. For $X$ uniform on $(0,1)$,
$\mathbb{E}[X^2]=1/3$, so that $\operatorname{Var}(X)=1/3-1/4=1/12\approx0.083333$.
Variance has squared units, whereas the standard deviation has the same units as $X$. One
costly warning: $\mathbb{E}[X^2]\ne\mathbb{E}[X]^2$, and the difference is precisely
the variance; for the uniform variable above, $0.25$ against $0.333333$.

\begin{resep}
Materials: a table of values $x_i$ and probabilities $p_X(x_i)$.
\begin{enumerate}
  \item Add up all the $p_X(x_i)$. If the result is not one, the table is not yet valid.
  \item Multiply each $x_i$ by $p_X(x_i)$, then add: that is $\mathbb{E}[X]$.
  \item Square each $x_i$ first, multiply by $p_X(x_i)$, then add:
        that is $\mathbb{E}[X^2]$.
  \item Subtract the square of the result of step 2 from the result of step 3, namely
        $\operatorname{Var}(X)=\mathbb{E}[X^2]-(\mathbb{E}[X])^2$; its square root
        is the standard deviation.
\end{enumerate}
Example: for $x=0,1,2,3$ with $p=0.1;\,0.2;\,0.3;\,0.4$, step 2 gives
$2$, step 3 gives $5$, and step 4 gives $1$.

Warning about the order: square $x_i$ before multiplying by the probability, and subtract
$\mu^2$ after $\mathbb{E}[X^2]$ has been fully added up. Swapping that order produces
a number that measures no width at all.
\end{resep}

\kamus{expectation}{harapan}{the balance point of a loaded ruler: one number that
marks where the values pile up}

\kamus{variance}{variansi}{a measure of how far typical values stray from the balance
point, computed from squared distances so that plus and minus signs do not cancel each other}

\section{Four Discrete Distributions}
\tingkat{intermediate}

Four families of discrete distributions suffice for most practical work: one
binary trial, the number of binary trials, the waiting time until the first success, and
rare events in a time interval.

\paragraph{Bernoulli.} $X$ Bernoulli with parameter $p\in[0,1]$ takes the value $1$ with
probability $p$ and $0$ with probability $1-p$, so that $p_X(x)=p^x(1-p)^{1-x}$ for
$x\in\{0,1\}$. From those two points, $\mathbb{E}[X]=p$ and
$\operatorname{Var}(X)=p(1-p)$; the variance is largest at $p=1/2$, namely $1/4$, exactly
when the two possibilities are balanced.

\paragraph{Binomial.} If $n$ independent Bernoulli trials with success probability $p$
are run, the number of successes $X$ is binomial with
\[
p_X(k)=\binom{n}{k}p^k(1-p)^{n-k},\qquad k=0,1,\dots,n .
\]
Its mean and variance are derived most easily by splitting $X$ into indicators,
not from that formula.

\begin{teorema}\index{binomial distribution}
If $X$ is binomial with parameters $n$ and $p$, then
$\mathbb{E}[X]=np$ and $\operatorname{Var}(X)=np(1-p)$.
\end{teorema}

\begin{proof}
Write $X=\sum_{i=1}^{n}X_i$ with $X_i=1$ if trial $i$ succeeds and $0$ if it
fails. Each $X_i$ is Bernoulli with parameter $p$, hence $\mathbb{E}[X_i]=p$ and
$\operatorname{Var}(X_i)=p(1-p)$. Linearity gives
$\mathbb{E}[X]=\sum_i\mathbb{E}[X_i]=np$ without using independence at all.

For the variance, independence is needed for the first time. Because the trials are independent, the pairs
$X_i,X_j$ with $i\ne j$ are independent too, so that
$\operatorname{Cov}(X_i,X_j)=p^2-p^2=0$. Hence
\[
\operatorname{Var}(X)=\sum_{i=1}^{n}\operatorname{Var}(X_i)
+2\sum_{i<j}\operatorname{Cov}(X_i,X_j)=np(1-p).
\]
\end{proof}

The variance formula is more fragile than the mean formula: the former requires independence,
the latter does not. For $n=20$ and $p=0.35$, $\mathbb{E}[X]=7$ and
$\operatorname{Var}(X)=4.55=91/20$, with standard deviation $2.133073$; the mass
around the center is $p_X(7)=0.184401$, $p_X(10)=0.068614$, $p_X(12)=0.013564$.
The tail is small but not zero: $P(X\ge12)=0.019579$ and $P(X\le4)=0.118197$, so the
rule ``reject a box if it contains twelve defects or more'' rejects about two percent of
boxes.

\paragraph{Geometric.} If independent Bernoulli trials are repeated until the first success,
the number of trials needed is $X$ geometric with
$p_X(k)=p(1-p)^{k-1}$ for $k=1,2,\dots$ Its total mass is one because of the geometric series:
$\sum_{k\ge1}p(1-p)^{k-1}=p/(1-(1-p))=1$. With $q=1-p$,
\[
\sum_{k\ge1}kq^{k-1}=\frac{d}{dq}\frac{1}{1-q}=\frac{1}{(1-q)^2}=\frac{1}{p^2},
\]
so that $\mathbb{E}[X]=p\cdot p^{-2}=1/p$ and, in a similar way,
$\operatorname{Var}(X)=(1-p)/p^2$ from $\mathbb{E}[X^2]=(2-p)/p^2$. For $p=0.2$:
$\mathbb{E}[X]=5$, $\operatorname{Var}(X)=20$, $P(X=5)=0.081920$ and
$P(X>10)=0.10737418$; a simulation of $10^5$ draws gives mean $4.99176$ with
standard error $0.01414$.

\paragraph{Poisson.} For $X$ Poisson with parameter $\lambda>0$,
\[
p_X(k)=e^{-\lambda}\frac{\lambda^k}{k!},\qquad k=0,1,2,\dots
\]
Its total mass is one because of the exponential series:
$e^{-\lambda}\sum_k\lambda^k/k!=1$. Its mean and variance are equal, and both
come out of a single trick. Because $k/k!=1/(k-1)!$,
\[
\mathbb{E}[X]=\sum_{k\ge1}k\,e^{-\lambda}\frac{\lambda^k}{k!}
=\lambda e^{-\lambda}\sum_{k\ge1}\frac{\lambda^{k-1}}{(k-1)!}=\lambda .
\]
Writing $k(k-1)/k!=1/(k-2)!$ gives $\mathbb{E}[X(X-1)]=\lambda^2$, hence
$\operatorname{Var}(X)=\lambda$. One parameter governs the center and the width at once, and that
restricts the Poisson: data whose variance far exceeds their mean (called
\emph{overdispersion}) do not fit the Poisson model.

For $\lambda=3$: $p_X(0)=0.049787$, $p_X(3)=0.224042$,
$P(X\ge8)=0.011905$. That last number answers the queueing question from the opening of the chapter:
with three requests per second, about $1.2$ percent of seconds contain eight
requests or more.

The Poisson is the limit of the binomial when $n$ is large, $p$ is small, and $np\to\lambda$. For
$n=1000$ and $p=0.003$, the binomial gives $P(0)=0.049563083$ whereas the Poisson
gives $e^{-3}=0.049787068$, a relative difference of $-4.5\times10^{-3}$. The error is small but
not zero and its sign changes with $k$; the Poisson is used as an approximation, not as an
identity.

\section{Three Continuous Distributions}
\tingkat{intermediate}

\paragraph{Uniform.} $X$ is uniform on $(a,b)$ if its density is $1/(b-a)$ inside
that interval and zero outside. Its distribution function rises linearly from $0$ to $1$, and
\[
\mathbb{E}[X]=\frac{a+b}{2},\qquad
\operatorname{Var}(X)=\frac{(b-a)^2}{12}.
\]
The factor $12$ comes from $\int_a^b(x-\tfrac{a+b}{2})^2\frac{dx}{b-a}=(b-a)^2/12$; for
$(0,1)$ the result is $1/12\approx0.083333$, the same as the computation through
$\mathbb{E}[X^2]$ in the previous section.

\paragraph{Exponential.} $X$ exponential with parameter $\beta>0$ has
density $f_X(x)=\beta^{-1}e^{-x/\beta}$ for $x>0$, with
$F_X(x)=1-e^{-x/\beta}$, $\mathbb{E}[X]=\beta$, and $\operatorname{Var}(X)=\beta^2$.
The parameter $\beta$ is the mean waiting time. The property that distinguishes it from every
other continuous distribution is the absence of memory.

\begin{proposisi}[Memoryless]\index{exponential distribution}
For $s,t\ge0$, $P(X>s+t\mid X>s)=P(X>t)$.
\end{proposisi}

\begin{proof}
$P(X>u)=e^{-u/\beta}$, so that
\[
P(X>s+t\mid X>s)=\frac{P(X>s+t)}{P(X>s)}
=\frac{e^{-(s+t)/\beta}}{e^{-s/\beta}}=e^{-t/\beta}=P(X>t).
\]
\end{proof}

A component that has already lasted five years is no more likely to last one more year
than a new component: sensible for a process without aging, absurd for a lamp that wears out.
For $\beta=2$, $P(X>3)=0.22313016$ and a simulation of $10^5$ draws
gives $0.22237$;
conditional on $X>2$, the empirical probability of $X>5$ is $0.22028$ against the theoretical
$0.22313$. The exponential is the continuous twin of the geometric: if arrivals follow a
Poisson process with rate $\lambda$, the gaps between arrivals are exponential with
$\beta=1/\lambda$.

\paragraph{Normal.} $X$ normal with parameters $\mu$ and $\sigma^2>0$ has
density
\[
f_X(x)=\frac{1}{\sigma\sqrt{2\pi}}
\exp\!\left(-\frac{(x-\mu)^2}{2\sigma^2}\right),\qquad x\in\mathbb{R}.
\]
The constant in front is not a choice: it is the only value that makes the area under the
curve equal to one. That the area is one is proved with polar coordinates in Chapter 15:
write $I^2$ as the double integral of $e^{-(x^2+y^2)/2}$, change to polar, and obtain
$I^2=2\pi$.

\begin{proposisi}
If $X\sim N(\mu,\sigma^2)$, then $Z=(X-\mu)/\sigma\sim N(0,1)$.
\end{proposisi}

\begin{proof}
$F_Z(z)=P(X\le\mu+\sigma z)=F_X(\mu+\sigma z)$. Differentiate with respect to $z$ by the chain
rule: $f_Z(z)=\sigma f_X(\mu+\sigma z)=\varphi(z)$, the standard normal density.
\end{proof}

The function $\Phi(z)=P(Z\le z)$ has no elementary form; it is computed numerically or
read from a table. Three numbers worth memorizing: $P(|Z|\le1)=0.682689492137$,
$P(|Z|\le2)=0.954499736104$, $P(|Z|\le3)=0.997300203937$. For a probability that is not
symmetric, subtract two values of $\Phi$: $P(1<Z<2)=0.135905122$. Its quantile,
$\Phi^{-1}(0.975)=1.959964$, appears in every 95 percent confidence interval in
Chapter 20.

The normal is not a more correct distribution than the others. It matters for two
different reasons: many natural quantities approach it, and the sum of many independent variables
tends to approach it whatever their original shape. The second reason is the central limit
theorem, proved in Chapter 19.

\kamus{normal distribution}{distribusi normal}{the symmetric bell curve: almost
all values pile up near the center, and the tails thin out very quickly}

\kamus{exponential distribution}{distribusi eksponensial}{waiting time in a process
without aging: one that has waited a long time is no closer to the next event}

\section{Joint Distributions, Covariance, and Correlation}
\tingkat{advanced}

Two separate distributions do not store information about whether the two quantities
tend to rise together. The tool for that is the joint distribution.

\begin{contoh}
Let $X,Y\in\{0,1\}$ with joint mass $P(0,0)=0.1$, $P(0,1)=0.2$,
$P(1,0)=0.3$, $P(1,1)=0.4$. Adding up rows and columns gives the marginals
$P(X=1)=0.7$ and $P(Y=1)=0.6$. If $X$ and $Y$ were independent, the mass at $(1,1)$
would have to be $0.42$, whereas the actual value is $0.4$; so the two are not independent,
even though the difference is small.
\end{contoh}

In the continuous case, the joint distribution is described by the joint density
$f_{X,Y}\ge0$ with $\iint f_{X,Y}=1$, and the marginals are obtained by
integrating one variable out: $f_X(x)=\int f_{X,Y}(x,y)\,dy$. Independence means
$f_{X,Y}(x,y)=f_X(x)f_Y(y)$ for all $x,y$.

\begin{definisi}[Covariance and correlation]\index{covariance}\index{correlation}
With $\mu_X=\mathbb{E}[X]$ and $\mu_Y=\mathbb{E}[Y]$,
\[
\operatorname{Cov}(X,Y)=\mathbb{E}\bigl[(X-\mu_X)(Y-\mu_Y)\bigr],
\qquad
\rho(X,Y)=\frac{\operatorname{Cov}(X,Y)}{\sigma_X\sigma_Y}.
\]
\end{definisi}

Positive covariance means the two variables tend to deviate in the same direction;
negative means in opposite directions. Correlation is the unitless version of covariance.

\begin{proposisi}
$\operatorname{Cov}(X,Y)=\mathbb{E}[XY]-\mathbb{E}[X]\mathbb{E}[Y]$ and
$-1\le\rho(X,Y)\le1$.
\end{proposisi}

\begin{proof}
Expand the product and use linearity. For the bound on correlation, consider the inner product
$\langle U,V\rangle=\mathbb{E}[UV]$ on centered variables; the Cauchy-Schwarz
inequality of Chapter 10 gives
$|\langle U,V\rangle|^2\le\langle U,U\rangle\langle V,V\rangle$, namely
$\operatorname{Cov}(X,Y)^2\le\sigma_X^2\sigma_Y^2$.
\end{proof}

\begin{teorema}
$\operatorname{Var}(X+Y)=\operatorname{Var}(X)+\operatorname{Var}(Y)
+2\operatorname{Cov}(X,Y)$.
\end{teorema}

\begin{proof}
Write $X+Y-\mu_X-\mu_Y=(X-\mu_X)+(Y-\mu_Y)$, square it, and then use linearity on
the three terms that appear; the cross term is exactly twice the covariance.
\end{proof}

Back to the table above: $\mathbb{E}[X]=7/10$, $\mathbb{E}[Y]=3/5$,
$\mathbb{E}[XY]=2/5$, so that $\operatorname{Cov}(X,Y)=-1/50=-0.02$. With
$\operatorname{Var}(X)=21/100$ and $\operatorname{Var}(Y)=6/25$, the correlation is
$-0.089087$ and $\operatorname{Var}(X+Y)=0.41$.

Zero covariance is often mistaken for independence, whereas the implication runs only one way:
independence implies zero covariance, not the reverse. Take $X$ taking the values $-1,0,1$ with
equal probability and $Y=X^2$; then $\mathbb{E}[X]=0$ and
$\mathbb{E}[XY]=\mathbb{E}[X^3]=0$, so that the correlation is zero, even though $Y$ is determined
completely by $X$. Covariance captures only dependence in the form of a straight line.

\kamus{covariance}{kovarians}{a measure of whether two quantities tend to rise and fall
together, computed from the product of their deviations}

\section{Transformations of Random Variables}
\tingkat{advanced}

This work comes up again and again: prices converted to logarithms, lengths to areas, waiting times
to rates. The question is always the same, namely how the distribution changes when
$Y=r(X)$.

In the discrete case the answer is simple: mass is moved from the old values to the new
values. Let $P(X=-1)=1/4$, $P(X=0)=1/2$, $P(X=1)=1/4$, and $Y=X^2$. Then
$P(Y=0)=1/2$ and $P(Y=1)=P(X=-1)+P(X=1)=1/2$; $Y$ has only two values, fewer
than $X$, because $r(x)=x^2$ is not one-to-one.

In the continuous case the mass is zero at every point, so what is moved is length,
and length changes scale. Three standard steps answer everything: determine
$A_y=\{x:r(x)\le y\}$, compute $F_Y(y)=\int_{A_y}f_X(x)\,dx$, then differentiate
$f_Y=F_Y'$.

\begin{teorema}[Monotone transformation]\index{transformation of random variables}
If $r$ is strictly monotone and differentiable, with inverse $s=r^{-1}$, then
\[
f_Y(y)=f_X\bigl(s(y)\bigr)\left|\frac{ds(y)}{dy}\right| .
\]
\end{teorema}

\begin{proof}
Suppose $r$ is strictly increasing. Then $r(x)\le y$ is equivalent to $x\le s(y)$, so that
$F_Y(y)=F_X(s(y))$, and the chain rule gives $f_Y(y)=f_X(s(y))s'(y)$ with
$s'(y)>0$. If $r$ is strictly decreasing, the inequality reverses: $F_Y(y)=1-F_X(s(y))$, so that
$f_Y(y)=-f_X(s(y))s'(y)=f_X(s(y))|s'(y)|$.
\end{proof}

The factor $|ds/dy|$ is local stretching: if $r$ stretches the scale, the density
decreases, because the same probability is spread over a longer interval. In more than
one dimension that factor becomes the Jacobian determinant, whose general form is in Chapter 15.
Figure~\ref{fig:transformasi-sumbu} shows that stretching for $Y=X^2$: four
intervals of equal length land on intervals that grow longer as $y$ approaches one,
while the density of $Y$ rises near zero, where the intervals shorten.

\begin{figure}[htbp]
\centering
\begin{tikzpicture}[font=\footnotesize,>=Stealth]
  \draw[->,layergrey] (-0.25,0) -- (7.3,0) node[right] {$x$};
  \foreach \px/\lab in {0/0,1.75/0.25,3.5/0.5,5.25/0.75,7/1}{
    \draw[layergrey] (\px,0.07) -- (\px,-0.07);
    \node[above,font=\scriptsize] at (\px,0.07) {$\lab$};
  }
  \node[font=\scriptsize,text=layergrey,anchor=west] at (0.15,0.55)
    {equal-length intervals};
  \draw[->,layergrey] (-0.25,-2.6) -- (7.3,-2.6) node[right] {$y=x^2$};
  \foreach \px/\lab/\anc in {0/0/east,0.44/0.0625/east,1.75/0.25/north,
                             3.94/0.5625/west,7/1/west}{
    \draw[layergrey] (\px,-2.53) -- (\px,-2.67);
    \node[below,font=\scriptsize,anchor=\anc] at (\px,-2.67) {$\lab$};
  }
  \foreach \px/\py in {0/0,1.75/0.44,3.5/1.75,5.25/3.94,7/7}{
    \draw[->,layerindigo] (\px,-0.05) -- (\py,-2.55);
  }
  \node[font=\scriptsize,text=layerindigo,align=center,anchor=west] at (1.0,-1.85)
    {interval shortens,\\ density rises};
  \node[font=\scriptsize,text=layerindigo,align=center,anchor=east] at (6.6,-1.85)
    {interval lengthens,\\ density falls};
  \begin{scope}[xshift=8.7cm,yshift=-2.6cm]
    \draw[->,layergrey] (0,0) -- (5.4,0) node[right] {$y$};
    \draw[->,layergrey] (0,0) -- (0,3.95) node[left] {$f_Y(y)$};
    \draw[layerplum,very thick,domain=0.12:1,samples=100,smooth,variable=\t]
      plot ({5.0*\t},{1.175/sqrt(\t)});
    \draw[layerplum,dashed] (0,0) -- (0,3.8);
    \node[below,font=\scriptsize,anchor=east] at (0,0) {$0$};
    \node[below,font=\scriptsize] at (1.25,0) {$0.25$};
    \node[below,font=\scriptsize,anchor=west] at (5.0,0) {$1$};
    \draw[layergrey,dotted] (5.0,0) -- (5.0,1.175);
    \node[font=\scriptsize,text=layerplum,anchor=west] at (0.25,3.45) {blows up near $0$};
  \end{scope}
\end{tikzpicture}
\caption{The transformation $Y=X^2$ for $X$ uniform on $(0,1)$, drawn on two axes.
The point $x=0.25$ lands at $y=0.0625$, $x=0.5$ at $y=0.25$, and $x=0.75$ at
$y=0.5625$; an interval of length $0.25$ near zero shortens to $0.0625$, whereas the same interval
near one lengthens to $0.4375$. The right-hand panel is the density
$f_Y(y)=1/(2\sqrt{y})$: its value is $2$ at $y=0.0625$, $1$ at $y=0.25$, and $0.5$ at
$y=1$. The density rises exactly where the interval shortens, so that the total area stays
one.}
\label{fig:transformasi-sumbu}
\end{figure}

If $r$ is not monotone, the theorem above does not apply directly, but the three standard
steps still work with $A_y$ a union of intervals.

\begin{contoh}
Let $X$ be uniform on $(-1,3)$, so that $f_X(x)=1/4$ inside it, and $Y=X^2$
taking values in $(0,9)$. For $0<y<1$ the set $A_y=[-\sqrt y,\sqrt y]$ lies entirely
inside $(-1,3)$, so $F_Y(y)=\tfrac14\cdot2\sqrt y=\sqrt y/2$; for $1\le y<9$
that set becomes $[-1,\sqrt y]$, so $F_Y(y)=(\sqrt y+1)/4$. Its derivative is
\[
f_Y(y)=\begin{cases}
\dfrac{1}{4\sqrt y}, & 0<y<1,\\[4pt]
\dfrac{1}{8\sqrt y}, & 1<y<9,\\[4pt]
0, & \text{elsewhere}.
\end{cases}
\]
The two branches add up to one, $\tfrac12$ and $\tfrac12$. The density of $Y$ blows up
near zero: $x^2$ compresses the intervals there.
\end{contoh}

If $X$ is uniform on $(0,1)$ and $Y=X^2$, the result is tidier: $F_Y(y)=\sqrt y$ and
$f_Y(y)=1/(2\sqrt y)$, which is checked numerically in the Lab. If $Y=-\ln X$, then
$F_Y(y)=1-e^{-y}$, exponential with $\beta=1$: taking the logarithm of a uniform number
produces an exponential waiting time, and the converse is what computers use.
If $U$ is uniform and $F$ is the distribution function we want, then $F^{-1}(U)$
has distribution $F$.

\begin{tantangan}
A set of lecture notes proposes the waiting-time density
\[
f(x)=\frac{1}{2\sqrt{x}},\qquad 0<x<4,
\]
on the grounds that it is ``positive, decreasing, and similar in shape to the example above''. From that density
the author computes $P(X\le4)=1$ and then concludes that the distribution is already valid.

Examine that claim. Determine the total area under $f$, repair the density in two different valid
ways, and explain why the repair is not unique.

Hint 1: what decides the matter is not the value of $f$ at a single point, but the whole region
under its curve.

Hint 2: compute $\int_0^4\tfrac12 x^{-1/2}\,dx$; the result is $2$, not $1$. For the
total to be one, the constant in front must be halved to $1/4$, or the interval must be
narrowed to $(0,1)$ with the constant unchanged.

If you get stuck, go back to the section ``Mass Functions and Density Functions''.
\end{tantangan}

\begin{untukanak}
Imagine a lamp on the floor and a ball between the lamp and the wall.
The shadow of the ball on the wall is much larger than the ball, and the farther the ball is from
the lamp, the larger the shadow. Nothing increases in amount: the paint on the ball
stays the same, it is merely spread over a larger area.

Now move the ball a little when it is already far from the lamp. Its shadow moves
a long way. When the ball is still close to the lamp, the same movement shifts the shadow only
a little. That is, one centimeter near the lamp becomes many centimeters on the wall,
whereas one centimeter far from the lamp becomes very many centimeters.

Changing a random variable works like that. Some regions are compressed and some
are stretched, and the resulting density has to adjust so that the total
stays one. If it does not adjust, probability will appear out of thin air or vanish
somewhere.
\end{untukanak}

\section{A Worked Example}
\tingkat{intermediate}

A production line fills boxes containing 20 fuses. Each fuse is defective with
probability $0.35$, independently of the others, and a box is rejected if it contains twelve defects or
more.

With $n=20$ and $p=0.35$, the rejection rule gives $P(X\ge12)=0.019579$ from
the binomial tail, not from a normal approximation. Even the most likely value, $X=7$, has
probability only $18.4$ percent, so that one box does not tell us much.

A simulation of $10^5$ boxes with \texttt{default\_rng(20261024)} gives mean
$7.005260$ and variance $4.552918$. The standard error of the mean is $0.006745$, so the
deviation is $+0.78$ standard errors; the standard error of the variance is about $0.020348$, so the
deviation is $+0.14$. Both are close, and that is to be expected: a simulation does not
prove a formula, it only fails to refute it.

The second check concerns the total probability mass of the Poisson distribution. The sum
$e^{-\lambda}\lambda^k/k!$ up to a sufficiently large $K$ must equal one. It is computed
in log space so that nothing overflows:

\begin{center}\small
\begin{tabular}{rrrr}
\hline
$\lambda$ & $K$ & sum & $1-\text{sum}$ \\
\hline
$0.5$ & 68 & $1.00000000000000000$ & $0$ \\
$1$ & 73 & $1.00000000000000000$ & $0$ \\
$3$ & 83 & $1.00000000000000022$ & $-2.2\times10^{-16}$ \\
$10$ & 107 & $1.00000000000000311$ & $-3.1\times10^{-15}$ \\
$50$ & 194 & $0.99999999999999778$ & $+2.2\times10^{-15}$ \\
$100$ & 280 & $1.00000000000003886$ & $-3.9\times10^{-14}$ \\
\hline
\end{tabular}
\end{center}

The largest deviation appears at $\lambda=100$, namely $3.9\times10^{-14}$, about
175 times the machine epsilon $2.2\times10^{-16}$. That is not a sign that the formula is wrong, but a
sign that adding $281$ terms of very different sizes loses a few of the last
digits; the first fourteen digits are still correct.

The third check concerns the total of the normal density, because the constant
$1/\sqrt{2\pi}$ is easily mistyped. Adaptive quadrature on $[-8,8]$ gives
$0.99999999999999878$, so the shortfall is $1.22\times10^{-15}$. Theory
predicts that shortfall as exactly the tail outside the interval, namely
$2\Phi(-8)=1.24\times10^{-15}$. The two numbers agree to two significant figures: what
is missing is indeed the tail, not an error in the constant. On the whole number line the result is
$0.99999999999999978$ with a printed error estimate of $1.0\times10^{-8}$, a conservative
estimate.

The fourth check, the transformation: the density $1/(2\sqrt y)$ has integral
$0.9999999999999992$, with a remainder of $8\times10^{-16}$ from the end at $0$ where the density
blows up. The simulation gives $F_Y(0.09)=0.29944$, $F_Y(0.25)=0.50003$,
$F_Y(0.49)=0.69944$, $F_Y(0.81)=0.90057$ against the theoretical $0.3$, $0.5$,
$0.7$, $0.9$. The standard error of each number is about $0.00145$ at $p=0.3$ and
$0.00095$ at $p=0.9$; all four deviations are below one standard error.

\kertas{Check by hand. Compute $\binom{20}{7}0.35^7 0.65^{13}$ and
compare it with $0.184401$; explain why $\sum_{k=0}^{20}p_X(k)=1$ follows from the
binomial theorem, not from measurement. Finally: the sample mean $7.005260$
differs from $7$ by $0.00526$. Is that a bug in the program? Answer with the
standard error $0.006745$.}

\section{Computational Lab}
\tingkat{advanced}

The following code is run with Python 3, NumPy 2.5.2, and SciPy 1.18.0. All the numbers
in the table and in the preceding paragraphs come from this run; the seed is fixed at
$20261024$, so that the results can be repeated.

\begin{lab}
\begin{verbatim}
import numpy as np
from math import exp, lgamma, log, sqrt, pi
from scipy import stats
from scipy.integrate import quad

rng = np.random.default_rng(20261024)
R = 100_000

# 1. binomial: exact formula vs Monte Carlo
n, p = 20, 0.35
mu, var = n * p, n * p * (1 - p)
sd = sqrt(var)
x = rng.binomial(n, p, size=R)
se = sd / sqrt(R)
print("Binomial n=20 p=0.35")
print(f"  theory E={mu:.6f} Var={var:.6f} sd={sd:.6f}")
print(f"  MC     mean={x.mean():.6f} Var={x.var(ddof=1):.6f}")
print(f"  SE mean={se:.6f} z={(x.mean()-mu)/se:+.4f}")
print(f"  SE Var  ={var*sqrt(2/(R-1)):.6f}")
print(f"  P(X>=12)={stats.binom.sf(11, n, p):.6f}")

# 2. Poisson pmf sum, computed in log space
print("\nPoisson pmf sum, k=0..K")
for lam in (0.5, 1.0, 3.0, 10.0, 50.0, 100.0):
    K = int(lam + 12 * sqrt(lam) + 60)
    s = sum(exp(-lam + k*log(lam) - lgamma(k+1)) for k in range(K+1))
    print(f"  lam={lam:6.1f} K={K:4d} sum={s:.17f}"
          f" 1-sum={1-s:+.2e}")

# 3. total normal density must be 1
phi = lambda t: exp(-t*t/2) / sqrt(2*pi)
print("\nStandard normal density")
print(f"  phi(0)={phi(0):.15f} 1/sqrt(2pi)={1/sqrt(2*pi):.15f}")
print(f"  quad [-8,8]  ={quad(phi, -8, 8)[0]:.17f}")
print(f"  remainder={1-quad(phi, -8, 8)[0]:.2e}"
      f"  2*Phi(-8)={2*stats.norm.cdf(-8):.2e}")
v, e = quad(phi, -np.inf, np.inf)
print(f"  quad (-inf,inf)={v:.17f} reported error={e:.1e}")

# 4. transformation X~Unif(0,1)
u = rng.uniform(0, 1, R)
y = u * u
print("\nY=X^2, X~Unif(0,1)")
print(f"  integral of f_Y={quad(lambda t: 1/(2*sqrt(t)), 0, 1)[0]:.16f}")
for yy in (0.09, 0.25, 0.49, 0.81):
    print(f"  F_Y({yy:.2f}) MC={np.mean(y<=yy):.5f} theory={sqrt(yy):.5f}")
z = -np.log(u)
print("Y=-ln(X), should be Exp(1)")
print(f"  mean={z.mean():.5f} SE={1/sqrt(R):.5f}"
      f" Var={z.var(ddof=1):.5f}")
print(f"  P(Y>2) MC={np.mean(z>2):.5f} theory={exp(-2):.5f}")

# 5. geometric and exponential
g = rng.geometric(0.2, R)
print("\nGeometric p=0.2: E=5 Var=20")
print(f"  mean={g.mean():.5f} SE={sqrt(20/R):.5f}"
      f" Var={g.var(ddof=1):.5f}")
w = rng.exponential(2.0, R)
print("Exponential beta=2: E=2 Var=4")
print(f"  mean={w.mean():.5f} SE={2/sqrt(R):.5f}"
      f" Var={w.var(ddof=1):.5f}")
print(f"  P(X>3) MC={np.mean(w>3):.5f} theory={exp(-1.5):.5f}")
print(f"  P(X>5|X>2) MC={np.mean(w[w>2]>5):.5f} theory={exp(-1.5):.5f}")
\end{verbatim}
\end{lab}

The output obtained, as it stands:

\begin{verbatim}
Binomial n=20 p=0.35
  theory E=7.000000 Var=4.550000 sd=2.133073
  MC     mean=7.005260 Var=4.552918
  SE mean=0.006745 z=+0.7798
  SE Var  =0.020348
  P(X>=12)=0.019579

Poisson pmf sum, k=0..K
  lam=   0.5 K=  68 sum=1.00000000000000000 1-sum=+0.00e+00
  lam=   1.0 K=  73 sum=1.00000000000000000 1-sum=+0.00e+00
  lam=   3.0 K=  83 sum=1.00000000000000022 1-sum=-2.22e-16
  lam=  10.0 K= 107 sum=1.00000000000000311 1-sum=-3.11e-15
  lam=  50.0 K= 194 sum=0.99999999999999778 1-sum=+2.22e-15
  lam= 100.0 K= 280 sum=1.00000000000003886 1-sum=-3.89e-14

Standard normal density
  phi(0)=0.398942280401433 1/sqrt(2pi)=0.398942280401433
  quad [-8,8]  =0.99999999999999878
  remainder=1.22e-15  2*Phi(-8)=1.24e-15
  quad (-inf,inf)=0.99999999999999978 reported error=1.0e-08

Y=X^2, X~Unif(0,1)
  integral of f_Y=0.9999999999999992
  F_Y(0.09) MC=0.29944 theory=0.30000
  F_Y(0.25) MC=0.50003 theory=0.50000
  F_Y(0.49) MC=0.69944 theory=0.70000
  F_Y(0.81) MC=0.90057 theory=0.90000
Y=-ln(X), should be Exp(1)
  mean=0.99811 SE=0.00316 Var=0.99623
  P(Y>2) MC=0.13478 theory=0.13534

Geometric p=0.2: E=5 Var=20
  mean=4.99176 SE=0.01414 Var=19.92571
Exponential beta=2: E=2 Var=4
  mean=1.99895 SE=0.00632 Var=3.98400
  P(X>3) MC=0.22237 theory=0.22313
  P(X>5|X>2) MC=0.22028 theory=0.22313
\end{verbatim}

Four things can be read from that output. First, every sample mean deviates by less
than one standard error: $z=+0.78$ for the binomial, whereas $0.99811$, $4.99176$,
and $1.99895$ against $1$, $5$, and $2$ for the other three cases. A deviation
of that size has to appear; if all of them were zero, something would be wrong.

Second, the sample variances are slightly below the theoretical value in the long-tailed case,
$19.92571$ against $20$, because a single very large draw dominates the sum
of squares. Third, the error of the Poisson summation grows with $\lambda$, from zero to
$10^{-14}$ as the number of terms increases from $74$ to $281$. Fourth, the resulting
density integrates to one to $16$ digits, and $F_Y$ matches $\sqrt y$ at four test
points; the second is more important, because an integral of one only checks the constant,
whereas agreement at four points checks the shape of the function.

That code does not prove a single theorem in this chapter; it checks that the formulas
written down give the numbers stated. The proof of $\mathbb{E}[X]=np$ is in the indicator
argument, the proof of $f_Y(y)=1/(2\sqrt y)$ in the computation of the distribution function;
the output of a program is data, not an argument.

\begin{untukdosen}
The order I suggest: two systems with the same mean but different tails, then
the definition of a random variable, then straight to the cumulative distribution function. Do one small
table on the board, draw the staircase, ask for $F(1.4)$, and only after that the density.
Resist the temptation to mention the integral before students see why summation
is no longer enough.

The two misconceptions that appear most often. First, $f_X(x)$ is read as a probability,
so that a density of value $5$ feels impossible and $P(X=x)$ is taken to equal
$f_X(x)$. Second, $P(X=x)=0$ is concluded to mean impossible, so that
$P(0\le X\le1)=1$ becomes a paradox. The answer is countable additivity: an uncountable
union is not added up.

Sign of readiness to move on: students can explain the difference between mass and density without
reading their notes, and can compute $\mathbb{E}[X^2]$ without finding the density of $X^2$.
\end{untukdosen}

\section{Summary}
\tingkat{basic}

\begin{ringkasan}
\begin{itemize}
\item A random variable is a function from the outcome space to the real numbers, not a random
      number. The distribution function $F_X(x)=P(X\le x)$ carries the whole information of the distribution:
      nondecreasing, right-continuous, with limits $0$ and $1$.
\item A mass function weights points and its values are at most one; a density function
      weights length, may exceed one, and $P(X=x)=0$ at every point. Probability
      zero is not impossibility, because only countable unions are added up.
\item Expectation is the balance point and is linear without any condition of independence;
      $\operatorname{Var}(X)=\mathbb{E}[X^2]-\mathbb{E}[X]^2$ is the mean squared
      deviation.
\item Binomial $\mathbb{E}=np$, $\operatorname{Var}=np(1-p)$, derived through
      indicators; geometric $\mathbb{E}=1/p$, $\operatorname{Var}=(1-p)/p^2$; Poisson
      $\mathbb{E}=\operatorname{Var}=\lambda$, the binomial limit for rare events.
\item The normal is standardized to $N(0,1)$ and read through $\Phi$; the exponential is
      memoryless, the continuous twin of the geometric.
\item Covariance measures only linear dependence: zero does not mean independent, and a
      monotone transformation changes the density by the factor $|ds/dy|$.
\end{itemize}
\end{ringkasan}

\section{Exercises}
\tingkat{intermediate}

\begin{latihan}
\begin{enumerate}
  \item Let $p_X(-2)=0.1$, $p_X(0)=0.4$, $p_X(3)=c$. Determine $c$, then
        compute $\mathbb{E}[X]$, $\mathbb{E}[X^2]$, and $\operatorname{Var}(X)$.
  \item Prove $P(a<X\le b)=F_X(b)-F_X(a)$, then explain why
        $P(a\le X\le b)$ can differ from $P(a<X\le b)$ for a discrete variable and
        why the two are equal for a continuous variable.
  \item Assess this claim: ``because $P(X=x)=0$ for every $x$, the event $X=x$
        is impossible.'' Name the part that is correct, the part that is wrong, and the role of countable
        additivity.
  \item Derive $\mathbb{E}[X^2]=(2-p)/p^2$ for the geometric by differentiating
        $\sum k^2q^{k-1}$ twice, then conclude $\operatorname{Var}(X)$ and check
        for $p=0.2$.
  \item Let $X\sim N(3,4)$. Compute $P(X>1)$, $P(1<X<5)$, and the 0.2 quantile of
        $X$; compare with \texttt{scipy.stats.norm}.
  \item Compute $\sum_{k=0}^{40}e^{-2.5}2.5^k/k!$ and $P(X\ge8)$ for the Poisson
        with $\lambda=2.5$, then compare with the binomial $n=1000$, $p=0.0025$;
        explain the direction of the difference.
  \item Prove that the exponential is memoryless, then explain why the geometric is also
        memoryless and why the uniform on the integers is not.
  \item For $X$ uniform on $(-1,3)$ and $Y=X^2$, derive $f_Y$, check
        $\int f_Y=1$, then compute $\mathbb{E}[Y]$ through $f_Y$ and through
        $\mathbb{E}[X^2]$.
  \item Construct two variables whose covariance is zero but which are not independent. Compute
        $\operatorname{Cov}$ and $\rho$ exactly, and exhibit a joint probability
        that proves the dependence.
  \item How many repetitions $R$ are needed for the standard error of the mean of a binomial with $n=20$, $p=0.35$
        to be at most $0.005$? And for the standard error of the variance estimate to be at
        most $0.02$?
\end{enumerate}

\noindent\textbf{Hints.} Item 3: separate the statement about probability from the
statement about the existence of a set. Item 4: differentiate $\sum_{k\ge0}q^k=1/(1-q)$
twice with respect to $q$, and remember that $\sum k^2q^{k-1}$ contains $\sum k(k-1)q^{k-1}$.
Item 8: for $0<y<1$ the set $A_y$ has two branches, for $y\ge1$ only one.
\end{latihan}

\begin{untukpeneliti}
What is genuinely proved in this chapter: the properties of the distribution function together with their proofs,
the relation of the density to the distribution function through the fundamental theorem of calculus, the formula for the
expectation of a function of a random variable in the discrete case, linearity of expectation, the variance formula,
the mean and variance of the binomial through indicators, the geometric and Poisson means through
series, memorylessness of the exponential, and the monotone transformation formula.

What is quoted, not proved: that a function satisfying the three properties of a distribution function
always comes from some probability space (Wasserman, Theorem 2.8, whose direction uses the
Caratheodory theorem), that the distribution function determines every probability (Wasserman,
Theorem 2.7), the normalization of the normal through polar coordinates (Chapter 15), and the correlation
bound through Cauchy-Schwarz (Chapter 10).

Primary sources: Wasserman, \emph{All of Statistics}, Sections 2.1--2.4 (including Lemma
2.15 and the warning that a density may exceed one) and 2.11 (Example 2.47 for
$Y=X^2$ on $U(-1,3)$), as well as 3.1--3.4. Blitzstein and Hwang, \emph{Introduction to
Probability}, Chapters 3--5 (LOTUS in 4.5, Poisson in 4.7, exponential in 5.5) and 7.3
for covariance.

Limits of validity: the density is not unique, so every statement about $f_X$
is a statement about the same class of functions almost everywhere; the Jacobian formula
requires monotonicity and differentiability; the Cauchy expectation does not exist, so every
``mean'' of very heavy-tailed data must be suspected first. Open question: how well a density can be estimated
from data without assuming a shape, and whether the tails of real distributions are thin enough to
justify the use of variance.
\end{untukpeneliti}

\peta{Chapter 18 follows Wasserman, \emph{All of Statistics}, Chapter 2 (2.1--2.4, Lemma 2.15,
and 2.11 for transformation) and Chapter 3 (3.1--3.4). In Blitzstein and Hwang,
\emph{Introduction to Probability}, Chapter 3 (3.6--3.7), Chapter 4 (4.2, 4.5, 4.7), Chapter 5
(5.1--5.5), and 7.3. The integrals used are the Riemann integrals of Chapter 13; change of
variables in more than one dimension is in Chapter 15.}
